% !TeX root = ../main.tex
%------------------------------------------------------------------------------
% Sustainability, Carbon and Resilience
%
% Enriched from the author's study report on sustainability and resilience of
% underground infrastructure, which supplies the module-based carbon
% workflow, the design-variable hierarchy and the resilience assessment
% sequence. Other sources: EN 15978; EN 17472; PAS 2080.
%------------------------------------------------------------------------------
\chapter{Sustainability, Carbon and Resilience}
\label{ch:sustainability}

Two lifecycle criteria have entered underground design over the past decade,
and this chapter treats them together because they are decided at the same
point and measured against the same reference. Carbon is determined by
decisions about alignment, diameter, lining thickness, support timing and
excavation method, which are also the decisions that govern safety and
serviceability; it is therefore a design quantity rather than a reporting
obligation. Resilience concerns the service the asset delivers over its life
and the rate at which that service is restored after disruption. Both are
assessed against the functionality of the asset rather than the capacity of
its components, and separating them encourages the mistaken view that
sustainability concerns materials while resilience concerns hazards. They are
better understood as two views of the same question: what this asset costs
and delivers across its life.

\section{Carbon sources and assessment}
\label{sec:carbon-sources}

\subsection{Where the carbon is}
\label{sec:embodied-carbon}

For a bored tunnel with a segmental lining, the concrete dominates
(Figure~\ref{fig:sus-carbon-elements}). Cement
clinker in the lining concrete is the single largest contributor. The order
of the remaining terms depends on the drive: reinforcement steel, annulus
grout and secondary concrete are set by the design, while the diesel and
electricity consumed by excavation, ventilation, dewatering and spoil haulage
depend on the length of the drive, the electricity grid and the disposal
site. In the case of Figure~\ref{fig:sus-carbon-elements}, machine
electricity from the grid exceeds the reinforcement, and spoil
haulage exceeds the grout. For a
sprayed concrete lined cavern the ranking shifts: shotcrete quantities are
harder to control, rebound losses count as material used, and the reinforcement
fraction is usually lower.

Two consequences follow. The first is that lining geometry (diameter and
thickness) is a first-order carbon variable, because it multiplies the
largest term directly. The second is that construction energy, often assumed
to be a minor contribution, becomes significant on long deep drives where
excavation, ventilation and cooling run continuously for years.

Spoil transport is the term most often underestimated. It scales with haulage
distance and vehicle type, both of which are set by disposal siting decisions
taken outside the engineering team, and it can rival the reinforcement
contribution on a project with no local disposal option.

\begin{figure}[htbp]
\centering
\datafigure[\textwidth]{ch03/sus-carbon-elements}
\caption{Carbon per metre of a bored rail tunnel by element, modules A1 to A5:
6.3~m internal diameter, 300~mm C50 segmental lining, machine on grid
electricity, 60~km spoil round trip; data from the baseline assessment of
\citet{chen2024carbonbaseline}. Lining concrete is the largest term; on this
drive the machine's electricity exceeds the reinforcement.}
\label{fig:sus-carbon-elements}
\end{figure}

\subsection{Life-cycle module boundaries}
\label{sec:carbon-modules}

Carbon figures are comparable only when the boundary is declared. The
EN~15978 and EN~17472 module structure~\citep{en15978,en17472} provides that declaration: modules
A1--A3 cover product manufacture, A4 transport to site, A5 construction
process, B use, maintenance and operational energy, C end of life, and D benefits beyond the
system boundary.

Most published tunnel carbon figures are not comparable, and the reason is
almost always boundary rather than method. A figure covering A1--A3 only will
be materially lower than one covering A1--A5 for the same design
(Figure~\ref{fig:sus-carbon-boundary}), by a margin
that widens with haulage distance and drive length, and a
whole-life figure including B and C may be higher again. A value quoted
without its modules is as incomplete as a number quoted without units.

\begin{figure}[htbp]
\centering
\datafigure[0.62\textwidth]{ch03/sus-carbon-boundary}
\caption{The tunnel of Figure~\ref{fig:sus-carbon-elements} reported under
three boundaries: 7.0~tCO$_2$e/m for A1--A3, 7.1 with transport and 10.6 with
the construction process. The A1--A3 figure is a third lower than the A1--A5
figure for an identical design.}
\label{fig:sus-carbon-boundary}
\end{figure}

The practical rule is that a carbon assessment must state four things
together: the quantities and their source; the carbon factors and their
source; the modules included; and the process assumptions behind the construction-stage
term. An assessment that omits any of the four cannot be audited, compared or
defended.

\subsection{Model-based quantity take-off and carbon factors}
\label{sec:model-based-carbon}

Quantities should come from the digital asset model of
Chapter~\ref{ch:digital-twin} (lining concrete, reinforcement, grout,
backfill, fit-out, spoil) rather than from a spreadsheet rebuilt at each
design stage. The benefit is consistency more than efficiency: when the quantity
model and the design model are the same object, a design change updates the
carbon estimate automatically, and the estimate stays current long enough to
influence the design.

Carbon factors require the same discipline. A factor is meaningful only with
its provenance: database or environmental product declaration, concrete grade
and supplementary cementitious material content, steel production route, and
the geography of the supply and its electricity grid. Two projects using the
same nominal concrete grade in different jurisdictions can differ by a
substantial margin on the product term alone.

\section{Carbon in design and construction}
\label{sec:carbon-in-design}

\subsection{The accounting logic}
\label{sec:carbon-formula}

The total for the product and construction stages assembles a product term
and a transport term, summed over the material inventory, and a project-level
process term:
\[
GWP_A \;=\; \sum_i V_i \left( GWP_{i,\mathrm{A1-A3}} + GWP_{i,\mathrm{A4}} \right)
\;+\; GWP_{\mathrm{A5}}
\]
where $GWP_A$ is the global warming potential of the product and construction
stages in
kilograms of carbon dioxide equivalent, $V_i$ is the quantity of material $i$
in the unit in which its factors are declared (cubic metres of concrete,
tonnes of steel), $GWP_{i,\mathrm{A1-A3}}$ and $GWP_{i,\mathrm{A4}}$ are its
product-stage and transport factors in kilograms of carbon dioxide equivalent
per unit of that quantity, so that any density conversion is carried in the
factor and not in $V_i$, and $GWP_{\mathrm{A5}}$ is the construction-process
term, also in kilograms of carbon dioxide equivalent; $i$ is here a summation
index. For mechanised excavation the
construction-stage term is dominated by machine energy, and is assembled by
energy carrier:
\[
GWP_{\mathrm{A5}} \;=\; \sum_j E_j\, EF_j ,
\qquad E_{\mathrm{TBM}} = \sum_j E_j ,
\]
where $E_j$ is the energy drawn as carrier $j$ (grid electricity and
diesel in most drives) in kilowatt-hours, diesel being converted from the
metered volume at its net calorific value, $EF_j$ is the emission factor of
that carrier, including its upstream supply, in kilograms of carbon dioxide
equivalent per kilowatt-hour, and $E_{\mathrm{TBM}}$ is the total energy
drawn by the excavation in kilowatt-hours.

\noindent\textbf{This is accounting logic, not a design equation.}\enspace
Neither expression predicts anything. They require every
reported figure to be traceable to a quantity, a boundary, a factor and a
stated modelling assumption. The excavation energy $E_{\mathrm{TBM}}$ in
particular should be estimated from ground--machine interaction (torque,
thrust, penetration rate, rotation speed, diameter, depth, ground class and
in-situ stress) using the performance methods of
Chapter~\ref{ch:tbm}, not assigned as a generic percentage of project
emissions.

\subsection{Carbon as a design variable}
\label{sec:carbon-hierarchy}

Carbon sensitivity is not evenly distributed across design decisions. In
approximate order of leverage:

\begin{enumerate}
 \item alignment, depth, diameter and geometric scale;
 \item lining thickness, concrete grade, supplementary cementitious material
 content and reinforcement ratio;
 \item support distance from the face, and the ground--support interaction it
 mobilises;
 \item excavation method and the disturbance it causes;
 \item machine performance (thrust, torque, penetration rate) and the
 construction energy that follows;
 \item spoil removal distance and transport mode;
 \item restoration strategy and its timing during operation.
\end{enumerate}

The ordering is the main design message of the chapter. Early geometric
decisions and construction control dominate the achievable reduction
(Figure~\ref{fig:sus-lining-carbon}); material
substitution is a late lever with a bounded ceiling. A project that defers
carbon to the specification stage has already lost most of the achievable
reduction.

\begin{figure}[htbp]
\centering
\datafigure[0.62\textwidth]{ch03/sus-lining-carbon}
\caption{Product-stage carbon of a segmental lining per metre for three ratios
of internal diameter to thickness: ring area times 811~kgCO$_2$e/m$^3$ (C50 at
598, plus 135~kg/m$^3$ of reinforcement at 1.58~kgCO$_2$e/kg). Carbon grows
with the square of diameter; $D_i/t = 18$ costs 40\,\% more than 25. Circle:
the tunnel of Figure~\ref{fig:sus-carbon-elements}.}
\label{fig:sus-lining-carbon}
\end{figure}

The third item is where
carbon depends directly on rock mechanics. Permitting more ground convergence before support
is mobilised reduces the support pressure the lining must carry, and therefore
its thickness and reinforcement. That is a carbon reduction obtained through
the ground reaction curve of Chapter~\ref{ch:ccm} rather than through a
material substitution. It is available only within serviceability, construction
and risk limits, and it holds only on the falling branch of
the ground reaction curve (Figure~\ref{fig:sus-support-timing}). In squeezing or strain-softening ground, loosening
beyond a critical convergence raises the pressure required, so that delaying
support increases the demand on it, and the lever has to be checked against
the curve for the ground in question.

\begin{figure}[htbp]
\centering
\datafigure[0.62\textwidth]{ch03/sus-support-timing}
\caption{Schematic of the support-timing lever in strain-softening ground.
The pressure required to hold the opening falls as more convergence is allowed
before support, reaches a minimum at a critical convergence, and rises beyond
it as the ground loosens. Delay saves lining, and carbon, only on the falling
branch.}
\label{fig:sus-support-timing}
\end{figure}

\subsection{Optimisation and trade-offs}
\label{sec:carbon-optimisation}

A defensible optimisation workflow defines the design variables; the
geotechnical inputs and their uncertainty ranges; the performance constraints
on factor of safety, support pressure, deformation, settlement,
constructability and durability; and the carbon objective by module. Design
alternatives are then evaluated using analytical, convergence--confinement,
numerical or surrogate models according to the project stage.

There is rarely a single minimum-carbon design. Multi-objective optimisation
yields a Pareto set within which no objective can be improved without
worsening another, and the preferred option is selected from that set on
stated grounds. Reporting the set rather than a single optimum shows
the client the trade-off that is being
accepted. The integrated treatment is Chapter~\ref{ch:integrated-design}.

\section{Materials, energy and environment}
\label{sec:materials-energy-environment}

\subsection{Low-carbon concrete and spoil}
\label{sec:materials-spoil}

Supplementary cementitious materials, alkali-activated binders and geopolymer
systems all reduce the product term, and all interact with the early-strength
requirement of tunnelling. A sprayed lining that must carry load
within hours, or a segment that must be demoulded on a production cycle,
cannot accept an arbitrary reduction in early strength development. The
achievable substitution is therefore bounded by the construction cycle, and it
is set jointly by the designer and the contractor rather than by the
specification alone.

Excavated material is the other large flow. Treated as waste it is a disposal
cost and an A5 carbon term; treated as a resource it can substitute for
quarried aggregate, fill and construction material. Reuse is
rarely limited by the rock itself; the usual limits are contamination, conditioning
residues from closed-face
excavation, and the regulatory classification of the material. Planning reuse
early (while the alignment, the excavation method and the disposal
strategy are still open) makes it achievable.

\subsection{Operational energy and thermal use}
\label{sec:operational-thermal}

Over a hundred-year life, ventilation, lighting, pumping and maintenance can
approach or exceed the embodied carbon of the structure. This matters for
design because some decisions trade one against the other. A drained lining
carries lower structural demand and less concrete, but commits the operator to
perpetual pumping and drainage maintenance; a tanked lining is heavier in
embodied terms and lighter in operation. Neither is universally correct, and
the comparison cannot be made without an operating assumption.

Underground structures also present an energy opportunity. Thermally activated
linings and absorber elements can recover geothermal heat from the ground, and
tunnels in operation reject substantial waste heat. The design implication is
cyclic thermal loading on the lining, which must be assessed rather than
assumed benign.

\subsection{Environmental approvals}
\label{sec:environmental-approvals}

Environmental assessment is where the sustainability argument becomes a
consent condition. Underground projects are typically assessed for groundwater
drawdown and its effect on dependent ecosystems and licensed bores, spoil
transport and disposal, construction noise and vibration, surface disturbance
at portals and shafts, and contaminated ground.

The engineering significance is that approval conditions become design
constraints, frequently binding ones. A drawdown limit set during approval
determines the lining drainage philosophy; a night-works restriction
determines the construction programme and therefore the ventilation and
haulage design. Approvals are best treated as an early input to design rather
than as a downstream permission to be obtained once the design is fixed.

Rating schemes provide the complementary framework, scoring a project against
defined sustainability criteria. They are useful for structuring an
improvement programme and for external comparability; they are not a
substitute for a quantified carbon assessment, and they measure process at
least as much as outcome.

\section{Resilience over the service life}
\label{sec:resilience}

\subsection{Functionality as the design measure}
\label{sec:functionality-measure}

A design that satisfies every capacity check can still leave an asset out of
service for months. A capacity check establishes whether the structure stands. It does not
establish whether the service continues, which is what the owner and the public
experience.

Resilience expresses performance as functionality over time. The asset begins
its life delivering a required level of service. Ageing, creep, corrosion and
leakage erode that level gradually; earthquakes, blasting, pressure transients
and construction events reduce it abruptly. Inspection, repair, strengthening,
grouting and relining restore it. Resilience is described by the shape of that curve, and
the useful metrics are drawn from it: the functionality lost immediately after
an event, the residual functionality, the delay before restoration begins, the
recovery trajectory, and the area under the curve over the recovery window.

Defining the service function is the first and most neglected step. It must be
expressed in terms of what the asset delivers (traffic capacity, flow
rate, safe occupancy, generating availability), with a stated threshold
below which the asset counts as out of service. Without that definition the
remaining metrics have no scale.

\subsection{Assessment and hazards}
\label{sec:resilience-assessment}

A workable assessment proceeds in eight steps: define the service function as
a functionality curve; establish baseline performance under ageing and
deterioration; apply the hazard scenarios; quantify the functionality lost
immediately after each event, giving vulnerability and robustness; estimate
residual functionality; estimate the delay before restoration and its
duration; evaluate the recovery trajectory; and compute a resilience index for
comparing intervention options. The index is the ratio of the area under the
functionality curve over the recovery window to the area that full service
would have delivered over the same window. It runs from 0 to 1, a higher
value being better, and it has to be quoted with its recovery window, on
which it depends (Figure~\ref{fig:sus-functionality}).

\begin{figure}[htbp]
\centering
\datafigure[\textwidth]{ch03/sus-functionality}
\caption{Two illustrative restorations after an event that leaves 35\,\% of
service, and the resilience index $R$ for each against the recovery window
$T$. A fast partial repair to 85\,\% scores higher over short windows; a
slower full restoration overtakes it beyond about 370~days. Shaded: the
service lost under the full restoration.}
\label{fig:sus-functionality}
\end{figure}

The hazards themselves are treated in their own chapters: seismic loading
in Chapter~\ref{ch:seismic}, squeezing in Chapter~\ref{ch:high-stress}, water inrush and gas in
Chapter~\ref{ch:ground-improvement}, deterioration in
Chapter~\ref{ch:assets}. What this framework adds is a common metric,
so that a fire, an earthquake and forty years of corrosion can be compared on
the same axis instead of being designed against independently.

Deterioration and discrete events also interact. A deteriorated lining is more
vulnerable to an earthquake, and an earthquake accelerates subsequent
deterioration. Treating the two as independent understates lifetime risk, and
the interaction is strongest where the consequence is highest, in
old assets carrying critical service.

\subsection{Restoration and recoverability}
\label{sec:restoration}

Restoration options are engineering design decisions with service and carbon
consequences, not maintenance activities to be selected on repair cost. The
comparison should be made on functionality recovered, time to recovery, cost,
constructability, risk reduction and carbon impact together. A cheaper repair
that takes three months longer may be the more expensive option once lost
service is counted.

Recoverability is itself designed in. Access for inspection and repair,
replaceable rather than monolithic elements, redundancy in critical systems,
and sacrificial detailing at fault crossings all determine how quickly an
asset can be brought back. These are decisions taken during design, usually
without an explicit resilience calculation to justify them, and they are among
the clearest cases where a lifecycle criterion changes a detail.

Uncertainty runs through the whole assessment (in hazard intensity, damage
state, detection delay and recovery rate), and reporting a single resilience
index without those bands overstates what the analysis establishes. The
methods for handling it are those of
Chapter~\ref{ch:integrated-design}. Finally, condition data and post-repair
performance should feed back into the model of
Chapter~\ref{ch:digital-twin}, updating the baseline for the next assessment
cycle. Without that return path the assessment remains a single study and is
not updated over the life of the asset.

\chapterreferences
